\section*{第 \ref{ch:b3:atmospheric-dynamics} 章——大气动力学}

\begin{solution}{exo:b3:atmospheric-dynamics:1}
(a)~$\zeta = 2\omega = 2\times2\pi/86\,400 = \qty{1.45e-4}{s^{-1}}$。
(b)~$f = \num{7.29e-5}$ 与 \qty{1.26e-4}{s^{-1}}：\omterm{def:b3:atmospheric-dynamics:eddy}{涡旋}的\omterm{def:b3:atmospheric-dynamics:vorticity-kinds}{行星涡度}约为
北纬 \ang{30} 处的两倍，略大于
北纬 \ang{60} 处的\omterm{def:b3:atmospheric-dynamics:vorticity-kinds}{行星涡度}。
\end{solution}

\begin{solution}{exo:b3:atmospheric-dynamics:2}
$\zeta = f_{30} - f_{40} = \num{7.29e-5} - \num{9.37e-5} =
\qty{-2.08e-5}{s^{-1}}$：反气旋式。
\end{solution}

\begin{solution}{exo:b3:atmospheric-dynamics:3}
$k = 2\pi/(5\times10^6) = \qty{1.257e-6}{m^{-1}}$；$c = -\beta/k^2 =
-\num{1.62e-11}/\num{1.58e-12} = \qty{-10.3}{m/s}$：向西。
\end{solution}

\begin{solution}{exo:b3:atmospheric-dynamics:4}
$\beta = 2\Omega\cos\ang{60}/a = \qty{1.14e-11}{m^{-1}.s^{-1}}$；
$\lambda_s = 2\pi\sqrt{10/\num{1.14e-11}} = \qty{5870}{km}$；圆周长
\qty{20015}{km}，可容纳 $n_s = 3.4$ 个波。
\end{solution}

\begin{solution}{exo:b3:atmospheric-dynamics:5}
$q = f/h$ 守恒：在山脉上空，$\zeta = f(8/10) - f =
\qty{-2.06e-5}{s^{-1}}$，为反气旋式；在盆地上空，$f(12/10) - f =
\qty{+2.06e-5}{s^{-1}}$，为气旋式。柱体被拉伸之处，即山脉背风一侧
与盆地上空，会产生气旋式涡度，
此处易于形成低压。
\end{solution}

\begin{solution}{exo:b3:atmospheric-dynamics:6}
$k = 5/(a\cos\ang{45}) = \qty{1.11e-6}{m^{-1}}$，$\beta/k^2 =
\qty{13.1}{m/s}$：$c = \qty{1.9}{m/s}$，$c_g = \qty{28.1}{m/s}$。能量
在 \num{2.1}~天内传播 \qty{5000}{km}，图样则需 \num{31}~天。
\end{solution}

\begin{solution}{exo:b3:atmospheric-dynamics:7}
(a)~$k^2 = \beta/U - l^2 = \num{1.079e-12} - \num{2.74e-13} =
\qty{8.05e-13}{m^{-2}}$，$k = \qty{8.97e-7}{m^{-1}}$：$\lambda =
\qty{7000}{km}$，$n = 4.0$，比 $l = 0$ 时波数更少。(b)~$l^2 =
\num{1.097e-12} > \beta/U = \num{1.079e-12}$：$k^2$ 将为负值。
\end{solution}

\begin{solution}{exo:b3:atmospheric-dynamics:8}
$L_d = 0.012\times10^4/\num{1.031e-4} = \qty{1164}{km}$；$\sigma =
0.31\times\num{1.031e-4}\times\num{2e-3}/0.012 = \qty{5.3e-6}{s^{-1}}$，即
$e$ 折时间为 \num{2.2}~天；波长 $2\pi\times1164/1.61 =
\qty{4540}{km}$。
\end{solution}

\begin{solution}{exo:b3:atmospheric-dynamics:9}
$[v^*T^*] = (VT_0/2)\cos\delta$：\qty{25}{K.m/s}、\qty{12.5}{K.m/s} 与 0；
热流量分别为 \num{25100}、\num{12600} 与 \qty{0}{W/m^2}。当
温度与风相位相差四分之一波长时，
向南输送的暖空气与向北的一样多：不存在输送。
\end{solution}

\begin{solution}{exo:b3:atmospheric-dynamics:10}
在 $U$ 附近线性化，取 $\zeta' = \nabla^2\psi'$ 与 $v' =
\partial_x\psi'$：$(\partial_t + U\partial_x)\nabla^2\psi' + \beta\partial_x
\psi' = 0$。设 $\psi' \propto \mathrm e^{\mathrm i(kx + ly - \omega t)}$，
则 $(\omega - Uk)K^2 = -\beta k$，故 $\omega = Uk - \beta k/K^2$。相对于
空气，任一波的相速为 $c - U = -\beta/K^2 < 0$。
\end{solution}

\begin{solution}{exo:b3:atmospheric-dynamics:11}
$\varphi_H = [5\times9.81\times12\,000\times0.25/(3\times(\num{7.292e-5})^2
\times(\num{6.371e6})^2)]^{1/2} = \qty{0.477}{rad}$，\ang{27.3}。
$u_M = \Omega a\varphi_H^2 = \qty{106}{m/s}$；精确地，$\Omega a\sin^2
\varphi_H/\cos\varphi_H = \qty{110}{m/s}$：为观测急流的
两倍半。
\end{solution}

\begin{solution}{exo:b3:atmospheric-dynamics:12}
(a)~在该带上，界面的位移为 $\eta = sy$，$-L/2 \le y
\le L/2$，相对于其水平位置。密度差减小的一层
受 $\eta$ 位移时，单位面积储存 $\tfrac12\rho g'\eta^2$ 的能量；
对整个带取平均，$\langle\eta^2\rangle = s^2L^2/12$。(b)~$g' = 0.327$：
\num{1.5e5} 与 \qty{6.1e5}{J/m^2}，相比之下，动能为 $\tfrac12\times0.7\times15^2
\times10^4 = \qty{7.9e5}{J/m^2}$：两者相当。
\end{solution}

\begin{solution}{pb:b3:atmospheric-dynamics:1}
\textbf{1.} $f_{45} = \qty{1.031e-4}{s^{-1}}$，$f_{50} =
\qty{1.117e-4}{s^{-1}}$；$\zeta = f_{45} - f_{50} = \qty{-8.6e-6}{s^{-1}}$。

\textbf{2.} $R = V/|\zeta| = 15/\num{8.6e-6} = \qty{1750}{km}$；顺时针，
向南回流。

\textbf{3.} $\zeta = f_{45}(9.5/10) - f_{50} = \qty{-1.38e-5}{s^{-1}}$，
反气旋性更强。

\textbf{4.} $\zeta = f_{45} - f_{40} = \qty{+9.4e-6}{s^{-1}}$，为气旋式：它
转而向北。无论向哪个方向偏离，气块都会获得
使其转回的涡度：这是一种恢复效应，并形成振荡。

\textbf{5.} $\beta = \qty{1.62e-11}{m^{-1}.s^{-1}}$；圆周长为
\qty{28306}{km}。

\textbf{6.} $n = 3$ 至 8：波长分别为 9435、7076、5661、4718、4044 与
\qty{3538}{km}；$c = -21.5$、$-5.5$、1.9、5.9、8.3 与 \qty{9.9}{m/s}；$c_g
= 51.5$、35.5、28.1、24.1、21.7 与 \qty{20.1}{m/s}。

\textbf{7.} $n = 3$ 与 4 向西移动，$n \ge 5$ 向东移动。波 $n = 5$ 需要
$10^6/1.86 = \qty{5.4e5}{s}$，即 \num{6.2}~天，方能移动 \qty{1000}{km}。

\textbf{8.} 下游，向东：当 $c_g = \qty{24.1}{m/s}$ 时，能量
在 \num{2.4}~天内传播 \qty{5000}{km}。

\textbf{9.} $\lambda_s = 2\pi\sqrt{15/\num{1.62e-11}} = \qty{6050}{km}$，
$n_s = 4.7$。

\textbf{10.} $n_s = 5.7$ 对应 \qty{10}{m/s}，4.1 对应 \qty{20}{m/s}：
较强的西风有利于更长的定常波，
绕半球的波数更少。

\textbf{11.} 在北纬 \ang{35}，$n_s = 5.8$；在北纬 \ang{55}，则为 3.4。

\textbf{12.} $\zeta = f(8/10) - f = \qty{-2.06e-5}{s^{-1}}$：气流经过山脉时受挤压，
使气柱呈反气旋性，背风侧的拉伸则使其呈气旋性。
山脉对气流施加了一种固定的涡度分布，
这是一种不移动的强迫。

\textbf{13.} 固定的强迫产生定常响应；响应
最强的情形是强迫图样本身为一个会静止不动的自由波，
即在 $\lambda_s$ 处发生共振。

\textbf{14.} $\Lambda = 25/10^4 = \qty{2.5e-3}{s^{-1}}$；
$|\partial T/\partial y| = \Lambda fT/g = \num{2.5e-3}\times\num{1.031e-4}
\times250/9.81 = \qty{6.6}{K}$，每 \qty{1000}{km}。

\textbf{15.} $\mathrm{Ri} = 0.011^2/(\num{2.5e-3})^2 = 19$：远大于 0.25，
不存在剪切湍流。

\textbf{16.} $L_d = 0.011\times10^4/\num{1.031e-4} = \qty{1067}{km}$。

\textbf{17.} $2\pi\times1067/1.61 = \qty{4160}{km}$；$28306/4160 = 6.8$
个波。

\textbf{18.} $2\pi\times1067/2.40 = \qty{2790}{km}$。

\textbf{19.} $\sigma = 0.31\times\num{1.031e-4}\times\num{2.5e-3}/0.011 =
\qty{7.3e-6}{s^{-1}}$，\num{0.63}~天$^{-1}$。

\textbf{20.} $\exp(3\times0.627) = 6.6$。

\textbf{21.} 风暴的波长 \qty{4160}{km}，短于定常波长
\qty{6050}{km}：作为罗斯贝波，它被西风向东输送，
而风暴则经由较慢或静止的长波
向东移动。

\textbf{22.} 随高度向西倾斜：只有这样，向极运动的空气
才比向赤道运动的空气更暖，热量向极地输送，
\omterm{def:b3:atmospheric-dynamics:ape}{有效位能}得以释放。

\textbf{23.} 增长最快的扰动每经过
\textbf{\num[text-series-to-math=true]{1.6}~天}，增长一个 $e$ 倍
（按计算为 \num{1.59}~天）。
\end{solution}
